\chapter{Sudden Vision Loss}

\section*{Definition}
Sudden vision loss is a rapid reduction or complete loss of sight in one or both eyes developing over seconds, minutes, or hours. It may be painless or associated with eye pain, headache, neurologic deficits, or flashes and floaters. It is a time-critical symptom because treatment may preserve vision and identify stroke or another emergency.

\section*{When to See a Doctor}
Call emergency services or obtain emergency ophthalmic assessment immediately for new complete or partial vision loss, a curtain or shadow across vision, sudden painless monocular blindness, severe eye pain with nausea, or vision loss with weakness, facial droop, speech difficulty, confusion, or severe headache. New vision loss with scalp tenderness, jaw pain, or fever is an emergency because giant-cell arteritis can permanently damage both eyes.

\section*{How It Develops}
Vision can fail when light cannot reach the retina, when the retina is damaged, when blood flow to the retina or optic nerve is interrupted, or when visual pathways in the brain are affected. Retinal artery occlusion, retinal detachment, optic-nerve ischemia, acute glaucoma, vitreous hemorrhage, and occipital stroke are important causes.

\section*{Causes}
\begin{itemize}
  \item Retinal detachment, retinal artery or vein occlusion, vitreous hemorrhage, and severe diabetic or hypertensive eye disease
  \item Ischemic optic neuropathy, optic neuritis, giant-cell arteritis, or compressive optic-nerve disease
  \item Acute angle-closure glaucoma, corneal disease, hyphema, or severe intraocular inflammation
  \item Stroke or transient ischemia affecting visual pathways; migraine aura usually causes temporary symptoms but requires assessment when new or atypical
  \item Trauma, medication or toxic exposure, and functional visual loss
\end{itemize}

\section*{Related Symptoms}
\begin{itemize}
  \item Flashes of light, new floaters, a curtain-like shadow, or distorted vision
  \item Eye pain, redness, headache, halos around lights, nausea, or vomiting
  \item Weakness, numbness, imbalance, speech difficulty, or loss of coordination
\end{itemize}

\section*{Medical Evaluation}
\begin{itemize}
  \item Visual acuity, pupils, visual fields, eye movements, color vision, intraocular pressure, and a dilated retinal examination are performed urgently.
  \item Optical coherence tomography, ocular ultrasound, or angiography may identify retinal or vitreous disease.
  \item Suspected stroke or vascular disease may require brain and vascular imaging, electrocardiography, and blood tests.
  \item In older adults with suspected giant-cell arteritis, inflammatory markers and urgent specialist evaluation are required.
\end{itemize}

\section*{Treatment Options}
Treatment depends on the cause and may include urgent retinal procedures, pressure-lowering therapy, corticosteroids for selected inflammatory or arteritic disease, treatment of vascular risk factors, or stroke care. Do not delay specialist treatment while trying home remedies.

\section*{Self-Care and Home Management}
Do not drive yourself, patch the eye, use leftover eye drops, or wait for symptoms to resolve. Note the exact time of onset, which eye is affected, associated neurologic symptoms, recent trauma, and medications. If a chemical entered the eye, rinse continuously with clean running water while arranging emergency care.

\section*{References}
\begin{thebibliography}{99}
\bibitem{sudden_vision_loss:ref1} American Academy of Ophthalmology. Retinal and Ophthalmic Artery Occlusions Preferred Practice Pattern. AAO; 2024.
\bibitem{sudden_vision_loss:ref2} American Academy of Ophthalmology. Preferred Practice Pattern: Retinal Detachment and Related Peripheral Vitreoretinal Disease. AAO; 2024.
\bibitem{sudden_vision_loss:ref3} Kleindorfer DO, et al. 2021 Guideline for the Prevention of Stroke in Patients With Stroke and Transient Ischemic Attack. \textit{Stroke}. 2021;52:e364--e467.
\end{thebibliography}
